\section{Acoustic transfer, radiation traction and temporal memory}\label{sec:acoustics}

\subsection{When a harmonic reduction is justified}
Let $\omega>0$ be angular carrier frequency, $T_a=2\pi/\omega$ its period,
$T_e$ an envelope variation time, $T_\Gamma$ an interface variation time, and
$\tau_r$ an acoustic energy ring-down time. A fast-time average over an interval
$\Delta t$ requires
\begin{equation}
 T_a\ll\Delta t\ll\min(T_e,T_\Gamma).
 \label{eq:averaging-separation}
\end{equation}
Replacing the field at each instant by a settled harmonic solution also requires
$\tau_r\ll\min(T_e,T_\Gamma)$. A short carrier period alone does not imply
short ring-down in a resonant cavity. When these conditions fail, causal wave
dynamics or the reference continuum equations must be retained.

For a fixed slow state, use the peak-phasor convention
$p_{1\ell}(\bm x,t)=\Rea\{P_\ell(\bm x)e^{-\mathrm i\omega t}\}$ and
$\bm v_{1\ell}(\bm x,t)=\Rea\{\bm V_\ell(\bm x)e^{-\mathrm i\omega t}\}$,
where $\mathrm i^2=-1$ and the subscript $1$ denotes first-order acoustics.
Thus the average of two real harmonic quantities with phasors $A,B$ is
$\frac12\Rea(AB^*)$. For locally homogeneous, inviscid, adiabatic fluids with
negligible mean advection, define the sound speed $c_\ell$ and obtain
\begin{equation}
 -\mathrm i\omega\rho_\ell\bm V_\ell=-\nabla P_\ell,\qquad
 -\mathrm i\omega P_\ell=-\rho_\ell c_\ell^2\nabla\cdot\bm V_\ell,
 \qquad
 \nabla\cdot\left(\frac{\nabla P_\ell}{\rho_\ell}\right)
       +\frac{\omega^2}{\rho_\ell c_\ell^2}P_\ell=0.
 \label{eq:linear-acoustics}
\end{equation}
The sound speed is the isentropic derivative of the equation of state. Loss,
thermal fields and elastic boundaries replace this simple operator when
needed; they are not unspecified tuning coefficients. Boundary-layer
reductions with declared scale conditions are developed in
\cite{bach2018,jorgensen2021}.

The acoustic interface conditions must follow by linearizing the kinematic and
traction conditions on the displaced interface. Let $\Xi$ be its first-order
normal-displacement phasor, in metres. About a quiescent interface,
$\bm V^-\cdot\nn=\bm V^+\cdot\nn=-\mathrm i\omega\Xi$. First-order traction
also contains the capillary curvature variation and the displacement through
base pressure gradients. Neglecting those corrections can yield continuity
of acoustic pressure and normal velocity, as in a common high-frequency
fluid--fluid scattering model \cite{chesneau2022}. That simplification requires
the fast interfacial restoring terms to be small; it is not an exact condition
for capillary waves at all frequencies. A compliant solid or a free surface
requires its own linearized boundary operator.

\subsection{The sign and content of radiation traction}
For the locally homogeneous inviscid acoustic reduction just specified, define
the time-averaged acoustic \emph{momentum flux}
\begin{equation}
 \bm R_\ell=
 \left(\frac{|P_\ell|^2}{4\rho_\ell c_\ell^2}
            -\frac{\rho_\ell|\bm V_\ell|^2}{4}\right)\Id
 +\frac{\rho_\ell}{2}\Rea(\bm V_\ell\otimes\bm V_\ell^*).
 \label{eq:radiation-flux}
\end{equation}
It is minus the radiation-stress convention used in
\cite{chesneau2022}; recording this sign is essential. The effective outward
interfacial traction and its scalar normal component are
\begin{equation}
 \bm t_{\rm ac}=(\bm R_- -\bm R_+)\nn=-\jump{\bm R}\nn,
 \qquad \Pi=\nn\cdot\bm t_{\rm ac}.
 \label{eq:acoustic-traction}
\end{equation}
Let $\bm T^h=-p^h\Id+\bm\tau(\bm U)$ denote slow hydrodynamic stress with the
mean-pressure convention associated with \cref{eq:radiation-flux}. The reduced
interface balance becomes
\begin{equation}
 \jump{\bm T^h}\nn+\bm t_{\rm ac}
       =\sigma\kappa\nn-\Gs\sigma.
 \label{eq:averaged-jump}
\end{equation}
In a consistent bulk momentum balance the acoustic force is
$-\nabla\cdot\bm R$ with the matching mean-pressure definition. A gradient
part can be absorbed into $p^h$; it must not be counted again at the interface.
In thermoviscous or inhomogeneous fluids, the averaged stress, heat source and
boundary-layer forcing have to be derived together from the reference
equations. A scalar normal pressure model does not cover all streaming or
tangential stresses.

Only in a one-sided pressure-release limit, with negligible adjoining-fluid
momentum flux, $P_-\simeq0$ and negligible tangential acoustic velocity at the
surface, does \cref{eq:acoustic-traction} reduce to
\begin{equation}
 \Pi\simeq\frac{\rho_-}{4}|\bm V_-\cdot\nn|^2.
 \label{eq:pressure-release-limit}
\end{equation}
This factor uses peak amplitudes. It must not be transferred unchanged to an
immersed liquid--liquid lens or to a surface supporting appreciable acoustic
pressure. Experiments on liquid interfaces and coupled calculations motivate
retaining the wave--shape interaction, rather than prescribing a load
independent of the deformed geometry \cite{bertin2012,chesneau2022,sisombat2023}.

\subsection{Impulses are causal histories, not instantaneous intensities}
For a fixed geometry and linear acoustics, let $\bm z_a(t)$ collect the wave
state, $\mathcal G_\Gamma(t)$ its causal propagation operator, and
$\mathcal B_\Gamma$ the command-to-source operator. Let $t_0$ be the starting
time and $\bm z_{a,0}$ the initial acoustic state. Then
\begin{equation}
 \bm z_a(t)=\mathcal G_\Gamma(t-t_0)\bm z_{a,0}
       +\int_{t_0}^t\mathcal G_\Gamma(t-s)
                         \mathcal B_\Gamma\bm a(s)\dd s.
 \label{eq:acoustic-memory}
\end{equation}
The radiative load is quadratic in this accumulated field. Two commands with
the same present amplitude can therefore yield different loads because their
histories differ. Changing the interface makes the operator time dependent,
so a fixed convolution kernel is then only a local approximation. Fluid
inertia, viscous relaxation, streaming, temperature and actuator mechanics
provide additional memory variables.

An actual impulse has a restricted mechanical interpretation. Consider a
declared inertial modal reduction with displacement coordinates $\bm q_f$,
modal mass matrix $\mathsf M_f$ and generalized acoustic force
$\bm f_a(t)$. Their component units are metres, kilograms and newtons,
respectively. Let $t_-,t_+$ bracket the effective force pulse, including any
relevant acoustic ring-down. If the shape and $\mathsf M_f$ change negligibly
over that interval and integrated restoring and dissipative forces are small,
integration of momentum gives
\begin{equation}
 \mathsf M_f[\dot{\bm q}_f(t_+)-\dot{\bm q}_f(t_-)]
       \simeq\int_{t_-}^{t_+}\bm f_a(t)\dd t.
 \label{eq:mechanical-impulse}
\end{equation}
The integral has units of momentum. In this limit a short force pulse primarily
changes velocity; the subsequent interface shape follows the inertial,
capillary and viscous evolution. This reduction does not apply when viscous
forces dominate the pulse interval, and an acoustic cycle average still needs
\cref{eq:averaging-separation}. Delivering the correct integrated force is not
the same as arriving at the target with negligible residual motion.

If several carriers $\omega_m$ are used, index $m$ labels their frequencies.
Quadratic products contain cross terms proportional to
$e^{-\mathrm i(\omega_m-\omega_n)t}$. They average away only with appropriate
coherence statistics or an averaging interval satisfying
$|\omega_m-\omega_n|\Delta t\gg1$, together with negligible relevant mechanical
response to the unresolved beats. Nearby carriers can drive slow beat forces;
a sum of independent intensities is not generally their physical load.

\subsection{A mean optical surface is not the entire optical state}
Even when $\Gamma$ denotes a stationary cycle-mean interface, the acoustic
displacement $\xi_a(\bm x,t)=\Rea(\Xi e^{-\mathrm i\omega t})$ can corrugate
its instantaneous position. Let $\delta n_o,\delta n_i$ denote acoustic or
thermal index perturbations along the incident and transmitted ray segments
$\mathcal L_o,\mathcal L_i$. For real conjugates and a nominal stationary
Fermat ray, the first variation of its optical path is
\begin{equation}
 \delta\Psi=(n_o\bm d_o-n_i\bm d_i)\cdot\nn\,\xi_a
       +\int_{\mathcal L_o}\delta n_o\dd s
       +\int_{\mathcal L_i}\delta n_i\dd s.
 \label{eq:fast-optical-error}
\end{equation}
Here $s$ is arclength, and a pupil-constant part is removed to obtain wavefront
error. Virtual conjugates use the corresponding signed path segments.
\Cref{eq:fast-optical-error} is a first-order estimate, not a full
acousto-optic scattering calculation. Material derivatives of optical index,
exposure time, and the fast interface response determine whether these terms
are negligible. A zero-mean wavefront fluctuation can still broaden an
exposure-averaged optical image. The eventual objective must therefore apply
to the intended optical operating state, not only to a plotted mean surface.
